% !TeX program = xelatex
% Original generic skeleton; replace its drafting text with computed content.
\documentclass[UTF8,a4paper,12pt]{ctexart}
\usepackage[left=2.5cm,right=2.5cm,top=2.5cm,bottom=2.5cm]{geometry}
\usepackage{amsmath,amssymb,bm,mathtools}
\usepackage{graphicx,float,booktabs,longtable,array,tabularx}
\usepackage{enumitem,caption,fancyhdr,url}
\usepackage[hidelinks,unicode]{hyperref}
\linespread{1.35}
\setlength{\parindent}{2em}
\setlength{\parskip}{0pt}
\renewcommand{\arraystretch}{1.38}
\setlist{leftmargin=2em,itemsep=0.2em}
\captionsetup{font=small,labelsep=quad}
\graphicspath{{figures/}}
\pagestyle{fancy}\fancyhf{}\fancyfoot[C]{\thepage}
\renewcommand{\headrulewidth}{0pt}
\ctexset{section={format=\centering\heiti\zihao{3},number=\chinese{section}、},subsection={format=\heiti\zihao{4}},subsubsection={format=\heiti\zihao{-4}}}
\newcommand{\dd}{\mathrm d}
\newcommand{\norm}[1]{\left\lVert#1\right\rVert}
\newcommand{\keywords}[1]{\par\medskip\noindent\textbf{关键词：}#1}
\newcommand{\paperfigure}[3]{\begin{figure}[!htbp]\centering\includegraphics[width=.86\linewidth,height=.38\textheight,keepaspectratio]{#1}\caption{#2}\label{#3}\end{figure}}
\title{数学建模论文}
\author{}\date{}
\begin{document}
\maketitle
\begin{abstract}
本文件为独立原创排版骨架。正式摘要在计算完成后，按总述、各问的问题与模型、关键数值和结论、简短总结组织，重要内容黑色加粗，并经实际编译调整至第一页。
\keywords{数学建模；数值计算；不确定性量化}
\end{abstract}
\clearpage
\tableofcontents
\clearpage
\section{简介}
\subsection{问题背景}
根据题目对象及已核验的相关研究组织背景。
\subsection{问题重述}
逐问写明输入、预测或优化任务及交付要求。
\section{问题分析}
由题目结构确定方法主线和逐问计算流程。
\section{模型假设}
列明进入方程、约束和数据处理的假设。
\section{名词解释及主要符号说明}
定义状态、参数、算子、指标及单位。
\section{模型建立与求解}
按各问给出数学推导、计算配置、数值结果与机制解释。
\section{灵敏度分析与鲁棒性检验}
比较同条件下的参数扰动与数据处理、边界或模型替换。
\section{模型的评价、改进与推广}
根据实际检验总结优点、误差来源、改进和推广。
\section*{AI工具使用声明}
根据实际使用情况及赛事规定填写。
\bibliographystyle{unsrt}
\bibliography{references}
\end{document}
